\documentclass[9pt, letterpaper]{article}

% Packages:
\usepackage[
    ignoreheadfoot, % set margins without considering header and footer
    top=0.3 cm, % seperation between body and page edge from the top
    bottom=0.3 cm, % seperation between body and page edge from the bottom
    left=0.7 cm, % seperation between body and page edge from the left
    right=0.7 cm, % seperation between body and page edge from the right
    footskip=0.2 cm, % seperation between body and footer
    % showframe % for debugging
]{geometry} % for adjusting page geometry
\usepackage{titlesec} % for customizing section titles
\usepackage{tabularx} % for making tables with fixed width columns
\usepackage{array} % tabularx requires this
\usepackage[dvipsnames]{xcolor} % for coloring text
\definecolor{primaryColor}{RGB}{0, 0, 0} % define primary color
\usepackage{enumitem} % for customizing lists
\usepackage{fontawesome5} % for using icons
\usepackage{amsmath} % for math
\usepackage[
    pdftitle={Resume},
    pdfauthor={Resume},
    pdfcreator={LaTeX with RenderCV},
    colorlinks=true,
    urlcolor=primaryColor
]{hyperref} % for links, metadata and bookmarks
\usepackage[pscoord]{eso-pic} % for floating text on the page
\usepackage{calc} % for calculating lengths
\usepackage{bookmark} % for bookmarks
\usepackage{lastpage} % for getting the total number of pages
\usepackage{changepage} % for one column entries (adjustwidth environment)
\usepackage{paracol} % for two and three column entries
\usepackage{ifthen} % for conditional statements
\usepackage{needspace} % for avoiding page brake right after the section title
\usepackage{iftex} % check if engine is pdflatex, xetex or luatex

% Ensure that generate pdf is machine readable/ATS parsable:
\ifPDFTeX
    \input{glyphtounicode}
    \pdfgentounicode=1
    \usepackage[T1]{fontenc}
    \usepackage[utf8]{inputenc}
    \usepackage{lmodern}
\fi

\usepackage{XCharter}

% Some settings:
\raggedright
\AtBeginEnvironment{adjustwidth}{\partopsep0pt} % remove space before adjustwidth environment
\pagestyle{empty} % no header or footer
\setcounter{secnumdepth}{0} % no section numbering
\setlength{\parindent}{0pt} % no indentation
\setlength{\topskip}{0pt} % no top skip
\setlength{\columnsep}{0.15cm} % set column seperation
\pagenumbering{gobble} % no page numbering

\titleformat{\section}{\needspace{4\baselineskip}\bfseries\large}{}{0pt}{}[\vspace{1pt}\titlerule]

\titlespacing{\section}{
    % left space:
    -1pt}{
    % top space:
    0.08 cm}{
    % bottom space:
    0.03 cm} % section title spacing

\renewcommand\labelitemi{$\vcenter{\hbox{\small$\bullet$}}$} % custom bullet points
\newenvironment{highlights}{
    \begin{itemize}[
        topsep=0.05 cm,
        parsep=0 pt,
        partopsep=0pt,
        itemsep=0pt,
        leftmargin=10pt
    ]
}{
    \end{itemize}
} % new environment for highlights


\newenvironment{highlightsforbulletentries}{
    \begin{itemize}[
        topsep=0.10 cm,
        parsep=0.10 cm,
        partopsep=0pt,
        itemsep=0pt,
        leftmargin=10pt
    ]
}{
    \end{itemize}
} % new environment for highlights for bullet entries

\newenvironment{onecolentry}{
    \begin{adjustwidth}{
        0 cm + 0.00001 cm}{
        0 cm + 0.00001 cm}
}{
    \end{adjustwidth}
} % new environment for one column entries

\newenvironment{twocolentry}[2][]{
    \onecolentry
    \def\secondColumn{#2}
    \setcolumnwidth{\fill, 4.5 cm}
    \begin{paracol}{2}
}{
    \switchcolumn \raggedleft \secondColumn
    \end{paracol}
    \endonecolentry
} % new environment for two column entries

\newenvironment{threecolentry}[3][]{
    \onecolentry
    \def\thirdColumn{#3}
    \setcolumnwidth{, \fill, 4.5 cm}
    \begin{paracol}{3}
    {\raggedright #2} \switchcolumn
}{
    \switchcolumn \raggedleft \thirdColumn
    \end{paracol}
    \endonecolentry
} % new environment for three column entries

\newenvironment{header}{
    \setlength{\topsep}{0pt}\par\kern\topsep\centering\linespread{1.5}
}{
    \par\kern\topsep
} % new environment for the header

\newcommand{\placelastupdatedtext}{% \placetextbox{<horizontal pos>}{<vertical pos>}{<stuff>}
  \AddToShipoutPictureFG*{% Add <stuff> to current page foreground
    \put(
        \LenToUnit{\paperwidth-2 cm-0 cm+0.05cm},
        \LenToUnit{\paperheight-1.0 cm}
    ){\vtop{{\null}\makebox[0pt][c]{
        \small\color{gray}\textit{Last updated in September 2025}\hspace{\widthof{Last updated in September 2025}}
    }}}%
  }%
}%

% save the original href command in a new command:
\let\hrefWithoutArrow\href

% new command for external links:


\begin{document}
    \newcommand{\AND}{\unskip
        \cleaders\copy\ANDbox\hskip\wd\ANDbox
        \ignorespaces
    }
    \newsavebox\ANDbox
    \sbox\ANDbox{$|$}

    \begin{header}
        \fontsize{16 pt}{26 pt}\selectfont Tanmay Bhuskute

        \normalsize
        \mbox{tanmay.v.bhuskute@gmail.com}%
        \kern 1.0 pt%
        \AND%
        \kern 1.0 pt%
        \mbox{linkedin.com/in/tanmay-bhuskute}%
        \kern 1.0 pt%
        \AND%
        \kern 1.0 pt%
        \mbox{github.com/Tanmayb05}%
        \kern 1.0 pt%
        \AND%
        \kern 1.0 pt%
        \mbox{+1 (732) 858-3027}%
        \kern 1.0 pt%
    \end{header}


\section{EDUCATION}

    \begin{twocolentry}{
        Aug 2024 - Dec 2026}
        \textbf{Master of Science in Computer Science} - University of Cincinnati \end{twocolentry}
    \begin{onecolentry}
    Relevant Coursework: Advanced Algorithms, Deep Learning, Intelligent Data Analysis, Database Theory.
    \end{onecolentry}
    \vspace{0.05 cm}

    \begin{twocolentry}{
        Aug 2018 - Jul 2022}
        \textbf{Bachelor of Engineering in Computer Engineering} - AISSMS College of Engineering \end{twocolentry}
    \begin{onecolentry}
    Relevant Coursework: Data Mining, Data Analytics, Machine Learning, Object-Oriented Programming.
    \end{onecolentry}

\section{TECHNICAL SKILLS}

    \begin{onecolentry}
    \textbf{Developer Tools:} AWS, Jenkins, Git, GitLab CI/CD, Docker, Kubernetes, ArgoCD, Helm, Terraform, Linux, Bash, SonarQube, Azure, Google Cloud Platform.
    \end{onecolentry}

    \begin{onecolentry}
    \textbf{Programming \& Scripting:} Python, Go, Java, SQL, JavaScript, TypeScript, C/C++.
    \end{onecolentry}

    \begin{onecolentry}
    \textbf{Web Backend:} REST APIs, FastAPI, Flask, Django, Node.js, React, HTML/CSS.
    \end{onecolentry}

    \begin{onecolentry}
    \textbf{Databases:} PostgreSQL, MySQL.
    \end{onecolentry}


\section{WORK EXPERIENCE}

    \begin{twocolentry}{Jul 2022 - Jun 2024}\textbf{Siemens Industry Software} (Software Development Engineer)\end{twocolentry}
    \begin{onecolentry}\textbf{Tech Stack:} AWS, Python, Boto3, REST APIs, Hashicorp Vault, Ansible, GitLab CI/CD, Shell, Kubernetes, ArgoCD, Helm, Amazon EFS \end{onecolentry}
    \begin{onecolentry}
        \begin{highlights}
            \item Automated Teamcenter X infrastructure workflows with \textbf{Python (Boto3)}, replacing legacy Ruby scripts and reducing deployment time by \textbf{40\%} across AWS environments
            \item Executed \textbf{44 production environment upgrades} in 4 months, diagnosing and resolving live infrastructure issues across Linux, EC2, databases, and Teamcenter X services
            \item Designed Ansible-based Blue-Green deployment framework, reducing production upgrade downtime by \textbf{50\%} while enabling reliable rollback and recovery
            \item Migrated 40+ production environments from process-based deployments to \textbf{Kubernetes} using ArgoCD, Helm, \textbf{GitLab CI/CD}, Amazon EFS, and HashiCorp Vault to modernize deployment architecture
            \item Served as Quality Coach and \textbf{SonarQube} Specialist for the Agile Release Train, increasing automated code coverage from 35\% to \textbf{78\%} by enforcing coding standards and CI/CD quality gates
            \item Mentored 5 software engineers through structured onboarding, technical documentation, and standardized development practices, accelerating team productivity
        \end{highlights}
    \end{onecolentry}
    \vspace{0.04 cm}

    \begin{twocolentry}{Mar 2022 - Jun 2022}\textbf{Siemens Industry Software} (Software Developer Intern)\end{twocolentry}
    \begin{onecolentry}\textbf{Tech Stack:} AWS, Terraform, Ansible, Terratest, Go \end{onecolentry}
    \begin{onecolentry}
        \begin{highlights}
            \item Automated provisioning of AWS resources with \textbf{Terraform}, reducing deployment time by \textbf{30\%}
            \item Developed Go + Terratest framework, catching \textbf{95\%} of misconfigurations pre-production
            \item Optimized DevOps workflows with \textbf{Ansible}, reducing environment setup time from hours to minutes
            \item Enhanced infrastructure validation with Terragrunt ensuring \textbf{100\%} successful Terraform runs
        \end{highlights}
    \end{onecolentry}

\section{PROJECTS AND RESEARCH}

    \begin{twocolentry}
        {May 2026}\textbf{PolicyPilot - Medical Benefit Policy Intelligence Portal}    \end{twocolentry}
    \begin{onecolentry}\textbf{Tech Stack:} React, TypeScript, Node.js, Express, Google Gemini, MCP, Docker, Zod, FHIR, Railway \end{onecolentry}
    \begin{onecolentry}
        \begin{highlights}
            \item Built Docker-based multi-stage deployment on Railway with health-check endpoints (\textbf{/api/health, /health}) and auto-restart on failure for a full-stack TypeScript healthcare app
            \item Designed 17-tool MCP server with \textbf{Zod schema validation} at every data boundary, eliminating a class of runtime type-mismatch bugs across all external tool calls
            \item Set up GitHub Actions CI/CD pipeline (test to build to deploy to health-check gate) with Vitest unit/integration tests and axe-core accessibility coverage
        \end{highlights}
    \end{onecolentry}
    \vspace{0.04 cm}

    \begin{twocolentry}
        {2026-07-05}\textbf{Spotify Personal Analytics Platform}    \end{twocolentry}
    \begin{onecolentry}\textbf{Tech Stack:} React 19.1.1, TypeScript 5.9.3, Vite 7.1.7, Material UI 7.3.4, MUI X Charts 8.14.0, Zustand 5.0.8, React Router 7.9.4, Axios 1.12.2, Python, FastAPI, Uvicorn, Pydantic, pandas, scikit-learn, PostgreSQL (Supabase), spotipy (Spotify Web API), lyricsgenius (Genius API), musixmatch-api, Netlify (frontend deploy), Render (backend deploy) \end{onecolentry}
    \begin{onecolentry}
        \begin{highlights}
            \item Engineered a 4-stage \textbf{Python ETL pipeline} processing 71K+ streaming events with custom rate limiting and exponential-backoff retry handling against a third-party API
            \item Implemented checkpointing and multi-source API fallback with automatic failure-based source disabling to keep long-running ingestion jobs resilient to mid-run failures
            \item Designed a \textbf{PostgreSQL} analytics schema with materialized views, composite/partial indexes, and 16 SQL RPC functions, enabling low-latency pre-aggregated queries
        \end{highlights}
    \end{onecolentry}

\section{EXTRACURRICULAR ACTIVITIES}

    \begin{onecolentry}
        \begin{highlights}
            \item Finalist at RevolutionUC 2025 Hackathon, among \textbf{300 participants from 31 schools}
            \item Earned \textbf{SAFe Practitioner 5.1 Certification} (Mar 2023)
        \end{highlights}
    \end{onecolentry}


\end{document}